% from ega0-21-fr.tex
\oldpage[0\textsubscript{IV-1}]{160}
par $q_\lambda$ l'homomorphisme canonique $E_\lambda\to C_\lambda$.
Comme par hypothèse $\mathfrak R_\lambda$ est nilpotent dans $E_\lambda$,
la première partie de la démonstration prouve l'existence et l'unicité d'un
$A$-homomorphisme $w_\lambda:B\to E_\lambda$ tel que
$p_\lambda\circ v=q_\lambda\circ w_\lambda$ et
$w_\lambda(x_\alpha)=z_{\alpha\lambda}$ pour tout $\alpha$.
L'unicité des $w_\lambda$ montre alors aussitôt que $(w_\lambda)$ est un
système projectif d'homomorphismes, et $w=\varprojlim w_\lambda$ répond
à la question.
\end{proof}

\begin{remark}[21.2.8]
\label{0.21.2.8-fr}
La vérification de l'existence et de l'unicité de l'homomorphisme $w:B\to E$
tel que $w(x_\alpha)=z_\alpha$ pour tout $\alpha$ peut se déduire directement
du fait que $B$ est une $A$-algèbre formellement lisse et que les
$d_{B/A}(x_\alpha)$ forment une base du $B$-module $\Omega^1_{B/A}$ (21.2.4),
sans faire intervenir le fait qu'il s'agit d'une $p$-base (de sorte que le
résultat est valable sans supposer les anneaux de caractéristique $p$).
En effet, on peut se borner au cas où $\mathfrak R^2=0$~; comme
$\operatorname{Der}_A(B,\mathfrak R)=\operatorname{Hom}_B(\Omega^1_{B/A},\mathfrak R)$
(20.4.8), il existe une $A$-dérivation $D$, et une seule, de $B$ dans
$\mathfrak R$ telle que $D(x_\alpha)=z'_\alpha$ pour toute famille
$(z'_\alpha)$ d'éléments de $\mathfrak R$~; la conclusion résulte de (20.1.1).
\end{remark}

\subsection{$p$-bases et modules d'imperfection}
\label{0.21.3-fr}

\begin{env}[21.3.1]
\label{0.21.3.1-fr}
Soient $A$, $B$ deux anneaux (de caractéristique $p$), $i:A\to B$,
$j:B\to A$ deux homomorphismes d'anneaux tels que
\begin{align}
\tag{21.3.1.1}\label{0.21.3.1.1-fr}
j(i(a))&=a^p\quad\text{pour tout }a\in A,\\
\tag{21.3.1.2}\label{0.21.3.1.2-fr}
i(j(b))&=b^p\quad\text{pour tout }b\in B.
\end{align}
Le plus souvent, $i$ sera injectif, de sorte que $A$ s'identifiera par
$i:A\to B$ à un sous-anneau de $B$~; une fois cette identification faite,
l'existence de $j$ implique que $B^p\subset A$, et $j$ s'identifie alors à $F_B$.
\end{env}

\begin{env}[21.3.2]
\label{0.21.3.2-fr}
Si $h:A^p\to A$ est l'injection canonique, on a, d'après (21.3.1.1) et
(21.3.1.2), un diagramme commutatif
\begin{equation}
\tag{21.3.2.1}\label{0.21.3.2.1-fr}
\xymatrix{
 B \ar[r]^{j} & A\\
 A \ar[u]^{i} \ar[r]^{F_A} & A^p \ar[u]_{h}
}
\end{equation}
si bien que le couple $(j,F_A)$ peut être considéré comme un di-homomorphisme
de la $A$-algèbre $B$ (pour $i$) dans la $A^p$-algèbre $A$ (pour $h$).
On en déduit un homomorphisme canonique de $A$-modules
\begin{equation}
\tag{21.3.2.2}\label{0.21.3.2.2-fr}
\pi_{B/A}:\Omega^1_{B/A}\otimes_B A_{[j]}\longrightarrow
\Omega^1_{A/A^p}=\Omega^1_A
\end{equation}
(cf. (20.5.4)~; l'identification de $\Omega^1_{A/A^p}$ et du module des
différentielles absolues $\Omega^1_A$ provient de (21.1.5), (ii)).
\begin{align}
\tag{21.3.2.3}\label{0.21.3.2.3-fr}
\Theta_{B/A}&=\Omega^1_{B/A}\otimes_B A_{[j]},\\
\tag{21.3.2.4}\label{0.21.3.2.4-fr}
\Xi_{B/A}&=\operatorname{Ker}(\pi_{B/A}).
\end{align}
